\chapter{Sites, Presheaves, and Locality}

Long before one speaks about hybrid dependent types, one needs a language for local
observation. The older sheaf-oriented work in the repository provides exactly that. A
program is not treated as one monolithic proof obligation. Instead it is covered by
observation sites: argument boundaries, branch guards, call results, error-sensitive
locations, output boundaries, and in the more expansive monograph, also SSA values,
wrapper boundaries, protocol or heap sites, proof sites, trace sites, and loop headers.

\begin{definition}[Program site category]
Given a program $P$, a program site category $\Site(P)$ is a category whose objects are
observation sites and whose morphisms are semantic transport maps induced by data flow,
control flow, wrapping, projection, or proof transport. A cover of a region of the
program is a family of sites whose joint information is intended to describe that region.
\end{definition}

The point of a site category is not category-theory theatrics. It is an implementation
stance. Once observations become sites, the analyzer can assign typed local sections to
those sites and restriction maps to the arrows between them. For example, an argument site
can carry a refinement about non-emptiness, a guard site can carry a branch-sensitive
fact, a call-result site can carry a transported summary from another function, and an
error site can carry a viability predicate describing when an operation is defined. The
question ``does the function type-check?'' is replaced by the more informative question
``do these local sections glue to a global contract?''

\begin{definition}[Semantic presheaf]
A semantic presheaf on a program site category is a contravariant functor
\[
  \Sem : \Site(P)^{\op} \to \mathsf{Lat}
\]
that assigns to each site a lattice of local sections and to each morphism a restriction
map transporting information backward along the corresponding semantic dependency.
\end{definition}

The repository repeatedly treats local sections as refinement-typed facts with enough
structure to be merged, compared, and explained. This is already visible in the older
paper material, where local sections are the atoms of bug finding and equivalence. It is
also visible in newer code, where bridges and checkers keep trying to harvest and move
site-local obligations.

\paragraph{Why the presheaf point of view matters.}
A presheaf gives \DepPy{} three important benefits.
\begin{enumerate}[leftmargin=*, itemsep=0.4em]
  \item It makes locality explicit. Every fact lives somewhere rather than vaguely on the
  whole program.
  \item It makes transport explicit. Facts are moved along named restrictions rather than by
  invisible global heuristics.
  \item It makes non-gluing diagnosable. Failures can be localized to overlaps and expressed
  as cocycles rather than as flat error lists.
\end{enumerate}

This is the conceptual hinge between the older sheaf-cohomology work and the newer hybrid
work. The hybrid system adds a layer dimension, but the idea of local sections and
restriction maps is already present at the foundational level.

\paragraph{Running example.}
Consider a simple normalization function.
\begin{lstlisting}
def normalize(values):
    total = sum(values)
    return [v / total for v in values]
\end{lstlisting}
A local section at the argument boundary may say ``values is a list of integers.'' A
local section at the division site may require ``$total \neq 0$.'' A local section at a call
site may transport assumptions from the caller. If those local sections do not agree on
the overlap, the failure should be represented as a mismatch in the presheaf rather than
as a mysterious one-line warning. That is the prototype of the cohomological story.

\begin{remark}
The repository sometimes writes as though the phrase ``bugs are $H^1$ classes'' is an
entire theory by itself. The more exact statement is that the site/presheaf framework
creates a mathematically principled place for gluing failures to live. The slogan is a
consequence of that framework, not a substitute for it.
\end{remark}
